\chapter{Физический приток и граница следовой аномалии}
\label{ch:v10-geometric-scale-beta-trace-anomaly-origin}

\section{Проверяемая гипотеза}

Первый гейт Тома X допустил растущую историю ячеек, но не вывел квантовый
ход относительно $\zeta=\log a$. Настоящий гейт проверяет, может ли
эндогенный приток новых физических ячеек одновременно нарушить точное
бозонно-духовое сокращение Тома IX, создать неустойчивость симметричного
вакуума и породить ненулевую интенсивную $\beta$-функцию.

Ключевое различие состоит в том, что калибровочные духи не должны перестать
сокращать фиктивные направления. Новый вклад обязан действовать только на
физический факторпространственный сектор и происходить от реальных входящих
степеней свободы.

\section{Закрытое сокращение}

Пусть одномодовый физический и духовый операторы совпадают:
\begin{equation}
 D_z(q)=D_{\rm gh}(q)=1+q^2.
 \label{eq:v10-inflow-closed-operators}
\end{equation}
Тогда
\begin{equation}
 Z_{\rm closed}(q)
 =\frac{\det D_{\rm gh}(q)}{\det D_z(q)}=1,
 \qquad
 \Gamma_{\rm closed}(q)=0.
 \label{eq:v10-inflow-closed-cancellation}
\end{equation}
Это сохраняет строгий запрет Тома IX: фиктивные калибровочные поля не
создают вакуумную кривизну.

\section{Нескомпенсированная физическая самоэнергия}

Положительный приток $J>0$ добавляет самоэнергию физической моде,
\begin{equation}
 D_z(q;J)=1+q^2+J,
 \qquad
 D_{\rm gh}(q)=1+q^2,
 \label{eq:v10-inflow-shifted-operator}
\end{equation}
но не меняет духовый оператор. Нормированный статистический множитель равен
\begin{equation}
 \widehat Z_J(q)
 =\frac{Z_J(q)}{Z_J(0)}
 =\frac{(1+J)(1+q^2)}{1+q^2+J}.
 \label{eq:v10-inflow-normalized-partition}
\end{equation}
Следовательно, физическое эффективное действие
\begin{equation}
 \Gamma_J(q)
 =-\log\widehat Z_J(q)
 =\log(1+q^2+J)-\log(1+q^2)-\log(1+J)
 \label{eq:v10-inflow-effective-action}
\end{equation}
не является плоским.

Его первые существенные вариации в симметричной точке имеют знаки
\begin{equation}
 \Gamma_J''(0)=-\frac{2J}{1+J}<0,
 \label{eq:v10-inflow-negative-curvature}
\end{equation}
\begin{equation}
 \Gamma_J^{(4)}(0)
 =\frac{12J(J+2)}{(J+1)^2}>0.
 \label{eq:v10-inflow-positive-fourth}
\end{equation}
Положительный физический приток поэтому создаёт отрицательную квадратичную
моду, отсутствовавшую в закрытом секторе.

\section{Точный свидетель нарушенной фазы}

Добавим уже допустимое положительное локальное насыщение $q^4/4$ и возьмём
рациональный свидетель $J=9/10$:
\begin{equation}
 V_{9/10}(q)=\frac{q^4}{4}+\Gamma_{9/10}(q).
 \label{eq:v10-inflow-witness-potential}
\end{equation}
Тогда
\begin{equation}
 V_{9/10}''(0)=-\frac{18}{19}<0,
 \label{eq:v10-inflow-witness-origin-hessian}
\end{equation}
а две ветви
\begin{equation}
 q_\star^\pm=\pm\frac{1}{\sqrt2}
 \label{eq:v10-inflow-broken-points}
\end{equation}
строго стационарны и имеют
\begin{equation}
 V_{9/10}''(q_\star^\pm)=\frac{37}{24}>0.
 \label{eq:v10-inflow-broken-hessian}
\end{equation}
Это точный конечномерный пример, в котором реальный приток создаёт
нарушенную фазу, тогда как закрытая бозонно-духовая система была плоской.
Значение $J=9/10$ является свидетелем существования архитектуры, а не
выведенным космологическим притоком.

\section{Экстенсивный рост и интенсивная граница}

Пусть число входящих степеней свободы масштабируется с пространственным
объёмом:
\begin{equation}
 J_{\rm ext}(\zeta)=J_0e^{3\zeta}.
 \label{eq:v10-inflow-extensive-current}
\end{equation}
Тогда формально
\begin{equation}
 \frac{dJ_{\rm ext}}{d\zeta}=3J_{\rm ext}.
 \label{eq:v10-inflow-extensive-beta}
\end{equation}
Однако это только инженерное масштабирование числа ячеек. После деления на
объёмный множитель получаем интенсивный приток
\begin{equation}
 j_{\rm int}(\zeta)
 :=e^{-3\zeta}J_{\rm ext}(\zeta)=J_0,
 \qquad
 \frac{dj_{\rm int}}{d\zeta}=0.
 \label{eq:v10-inflow-intensive-beta-zero}
\end{equation}
Следовательно, чистое умножение одинаковых ячеек даёт ненулевой
экстенсивный показатель, но не создаёт квантовую интенсивную
$\beta$-функцию или следовую аномалию.

\begin{theorem}[Приток нарушает сокращение в физическом секторе]
Положительная самоэнергия $J$, действующая на физическую моду и не
действующая на дух, создаёт ненулевой нормированный определитель,
отрицательную квадратичную вариацию и положительную четвёртую вариацию.
В присутствии положительного квартичного насыщения существует точный
локально устойчивый нарушенный свидетель $q=\pm1/\sqrt2$.
\end{theorem}

\begin{nogocriterion}[Геометрический объём не является квантовой аномалией]
Закон $J_{\rm ext}\propto e^{3\zeta}$ отражает только рост числа одинаковых
ячеек. После перехода к интенсивной величине ход равен нулю. Поэтому
ненулевая физическая $\beta$-функция требует зависимости локальной
самоэнергии, спектра или взаимодействия от $\zeta$; один объёмный приток её
не выводит.
\end{nogocriterion}

\begin{protocol}[Статус гейта притока]\begin{enumerate}
\item закрытое бозонно-духовое сокращение: \textbf{$1$};
\item нескомпенсированный физический определитель: \textbf{получен};
\item архитектура нарушения симметрии: \textbf{$6/6$};
\item точный нарушенный свидетель: \textbf{$q=\pm1/\sqrt2$};
\item интенсивная квантовая $\beta$-функция: \textbf{$0/1$};
\item физическое происхождение резервуара и аномалии: \textbf{$0/2$};
\item следующий гейт: \textbf{масштабно-зависимая спектральная самоэнергия
притока}.
\end{enumerate}\end{protocol}

Машинный сертификат сохранён в
\path{s2t_v10_geometric_scale_beta_trace_anomaly_origin_gate_results.json}.

\begin{thebibliography}{9}
\bibitem{v10-inflow-keldysh}
L.~M.~Sieberer, M.~Buchhold, S.~Diehl,
\emph{Keldysh Field Theory for Driven Open Quantum Systems},
Reports on Progress in Physics 79 (2016), 096001.

\bibitem{v10-inflow-repeated-interaction}
S.~Attal, Y.~Pautrat,
\emph{From Repeated to Continuous Quantum Interactions},
Annales Henri Poincar\'e 7 (2006), 59--104.
\end{thebibliography}